\documentclass[10pt]{extarticle}
\usepackage[a4paper,margin=1.5cm]{geometry}
\usepackage[T1]{fontenc}
\usepackage{amsmath,amssymb}
\usepackage{graphicx}
\usepackage[export]{adjustbox}
\usepackage{multicol}
\usepackage{enumitem}
\usepackage{booktabs,array,tabularx}
\setlength{\columnsep}{0.7cm}
\setlength{\parindent}{0pt}
\setlength{\parskip}{3pt}
\pagestyle{empty}
\makeatletter
% figure registry (figdefs.tex is generated by the extract step)
\newcommand{\qfigdef}[4]{\@namedef{qf@file@#1}{#2}\@namedef{qf@w@#1}{#3}\@namedef{qf@wide@#1}{#4}}
\newcommand{\qf@check}[1]{\@ifundefined{qf@file@#1}{\PackageError{score2}{Unknown figure id #1}{}}{}}
\newcommand{\qf@img}[4][]{\includegraphics[width=#2\dimexpr\@nameuse{qf@w@#4}pt\relax,max width=#3,max height=0.45\textheight,#1]{\@nameuse{qf@file@#4}}}
% Full-width questions: finalize wraps every question that contains a wide figure in \qfullbegin...\qfullend,
% so the WHOLE question (stem, figures, options) is set across the page in reading order, then the two columns
% resume. The question's head - its number, stem and first figure (\qfullsplit marks the end of the head) -
% is kept on one page, so a number is never left at a page foot with its figure overleaf; the rest may break
% across pages normally, so no big gaps are left. Without a \qfullsplit, the whole question is kept together
% if it is under 45% of a page. Never write these yourself.
\newif\ifqf@full
\newif\ifqf@boxed
\newsavebox{\qf@box}
\newcommand{\qf@open}{\global\qf@fulltrue\par\global\qf@boxedtrue
  \setbox\qf@box\vbox\bgroup\hsize\textwidth\linewidth\textwidth}
\newcommand{\qf@emit}[1]{\par\egroup\global\qf@boxedfalse
  \ifdim\dimexpr\ht\qf@box+\dp\qf@box\relax<#1\textheight\noindent\usebox{\qf@box}\par\else\unvbox\qf@box\fi}
\newcommand{\qfullsplit}{\ifqf@boxed\qf@emit{0.85}\fi}
\newcommand{\qf@close}{\ifqf@boxed\qf@emit{0.45}\fi\global\qf@fullfalse}
\newcommand{\qfullbegin}{\end{multicols}\qf@open}
\newcommand{\qfullend}{\qf@close\begin{multicols}{2}}
\newcommand{\qfullstart}{\qf@open}
\newcommand{\qfulldone}{\qf@close}
% \qfig[scale]{id}: block figure, centred, directly where it is written (size set by the tool).
% Fallback if a wide figure is not inside a full-width question: break out of the columns in place - no floats.
\newcommand{\qfig}[2][1]{\qf@check{#2}%
  \ifqf@full\qf@blk{#1}{#2}\else
  \ifnum\@nameuse{qf@wide@#2}=1 \expandafter\@firstoftwo\else\expandafter\@secondoftwo\fi
  {\end{multicols}\begin{center}\qf@img{#1}{\textwidth}{#2}\end{center}\begin{multicols}{2}}%
  {\qf@blk{#1}{#2}}\fi}
\newcommand{\qf@blk}[2]{\par\smallskip{\centering\qf@img{#1}{\linewidth}{#2}\par}\smallskip}
% \qfigi[scale]{id}: inline figure (inside option lists / tables / minipages); never breaks columns.
\newcommand{\qfigi}[2][1]{\qf@check{#2}\qf@img[valign=c]{#1}{\linewidth}{#2}}
\makeatother
% structure macros - lint reads these to check completeness and figure placement
\newcommand{\papersec}[1]{\par\bigskip{\centering\Large\bfseries PAPER-#1\par}\medskip}
\newcommand{\partsec}[2]{\par\medskip\noindent{\large\bfseries PART-#1: #2}\par\smallskip}
\newcommand{\sect}[1]{\par\smallskip\noindent\textbf{#1}\par}
\newcommand{\qn}[1]{\par\smallskip\noindent\textbf{#1.}~}
\newcommand{\paraq}[2]{\par\smallskip\noindent\textbf{Paragraph for Questions #1 and #2}\par}
\newenvironment{qoptions}{\begin{itemize}[leftmargin=1.8em,itemsep=1pt,topsep=2pt,label={}]}{\end{itemize}}
\newcommand{\opt}[2]{\item[(#1)] #2}
% figure registry (generated by score2_tools.py finalize)
\qfigdef{p1_phy_s1_q01_1}{TOAS - Test 14 (07 Aug 2022)-img/p1_phy_s1_q01_1.png}{84.6}{0}
\qfigdef{p1_phy_s1_q02_1}{TOAS - Test 14 (07 Aug 2022)-img/p1_phy_s1_q02_1.png}{103.2}{0}
\qfigdef{p1_phy_s1_q03_1}{TOAS - Test 14 (07 Aug 2022)-img/p1_phy_s1_q03_1.png}{100.4}{0}
\qfigdef{p1_phy_s1_q04_1}{TOAS - Test 14 (07 Aug 2022)-img/p1_phy_s1_q04_1.png}{79.4}{0}
\qfigdef{p1_phy_s1_q05_1}{TOAS - Test 14 (07 Aug 2022)-img/p1_phy_s1_q05_1.png}{140.1}{0}
\qfigdef{p1_phy_s1_q06_1}{TOAS - Test 14 (07 Aug 2022)-img/p1_phy_s1_q06_1.png}{171.7}{0}
\qfigdef{p1_phy_s1_q07_1}{TOAS - Test 14 (07 Aug 2022)-img/p1_phy_s1_q07_1.png}{134.7}{0}
\qfigdef{p1_phy_s1_q07_A}{TOAS - Test 14 (07 Aug 2022)-img/p1_phy_s1_q07_A.png}{237.0}{0}
\qfigdef{p1_phy_s1_q07_B}{TOAS - Test 14 (07 Aug 2022)-img/p1_phy_s1_q07_B.png}{227.5}{0}
\qfigdef{p1_phy_s1_q07_C}{TOAS - Test 14 (07 Aug 2022)-img/p1_phy_s1_q07_C.png}{277.0}{1}
\qfigdef{p1_phy_s1_q07_D}{TOAS - Test 14 (07 Aug 2022)-img/p1_phy_s1_q07_D.png}{133.9}{0}
\qfigdef{p1_phy_s1_q08_1}{TOAS - Test 14 (07 Aug 2022)-img/p1_phy_s1_q08_1.png}{134.7}{0}
\qfigdef{p1_phy_s2_para1_2_1}{TOAS - Test 14 (07 Aug 2022)-img/p1_phy_s2_para1_2_1.png}{132.2}{0}
\qfigdef{p1_phy_s2_para3_4_1}{TOAS - Test 14 (07 Aug 2022)-img/p1_phy_s2_para3_4_1.png}{134.7}{0}
\qfigdef{p1_phy_s2_para5_6_1}{TOAS - Test 14 (07 Aug 2022)-img/p1_phy_s2_para5_6_1.png}{92.6}{0}
\qfigdef{p1_phy_s3_q01_1}{TOAS - Test 14 (07 Aug 2022)-img/p1_phy_s3_q01_1.png}{163.8}{0}
\qfigdef{p1_phy_s3_q04_1}{TOAS - Test 14 (07 Aug 2022)-img/p1_phy_s3_q04_1.png}{134.7}{0}
\qfigdef{p1_phy_s3_q06_1}{TOAS - Test 14 (07 Aug 2022)-img/p1_phy_s3_q06_1.png}{134.4}{0}
\qfigdef{p1_chem_s1_q05_A}{TOAS - Test 14 (07 Aug 2022)-img/p1_chem_s1_q05_A.png}{55.8}{0}
\qfigdef{p1_chem_s1_q05_B}{TOAS - Test 14 (07 Aug 2022)-img/p1_chem_s1_q05_B.png}{66.3}{0}
\qfigdef{p1_chem_s1_q05_C}{TOAS - Test 14 (07 Aug 2022)-img/p1_chem_s1_q05_C.png}{61.0}{0}
\qfigdef{p1_chem_s1_q05_D}{TOAS - Test 14 (07 Aug 2022)-img/p1_chem_s1_q05_D.png}{61.0}{0}
\qfigdef{p1_chem_s1_q07_1}{TOAS - Test 14 (07 Aug 2022)-img/p1_chem_s1_q07_1.png}{138.9}{0}
\qfigdef{p1_chem_s2_q02_1}{TOAS - Test 14 (07 Aug 2022)-img/p1_chem_s2_q02_1.png}{239.0}{0}
\qfigdef{p1_chem_s2_para5_6_2}{TOAS - Test 14 (07 Aug 2022)-img/p1_chem_s2_para5_6_2.png}{141.1}{0}
\qfigdef{p1_chem_s3_q02_1}{TOAS - Test 14 (07 Aug 2022)-img/p1_chem_s3_q02_1.png}{102.1}{0}
\qfigdef{p1_chem_s3_q06_1}{TOAS - Test 14 (07 Aug 2022)-img/p1_chem_s3_q06_1.png}{245.2}{0}
\qfigdef{p1_chem_s3_q06_2}{TOAS - Test 14 (07 Aug 2022)-img/p1_chem_s3_q06_2.png}{245.2}{0}
\qfigdef{p1_chem_s3_q06_3}{TOAS - Test 14 (07 Aug 2022)-img/p1_chem_s3_q06_3.png}{268.7}{1}
\qfigdef{p1_chem_s3_q06_4}{TOAS - Test 14 (07 Aug 2022)-img/p1_chem_s3_q06_4.png}{307.3}{1}
\qfigdef{p1_chem_s3_q06_5}{TOAS - Test 14 (07 Aug 2022)-img/p1_chem_s3_q06_5.png}{245.2}{0}
\raggedcolumns
\makeatletter\@beginparpenalty=10000\makeatother
\newcolumntype{L}[1]{>{\raggedright\arraybackslash}p{#1}}
\begin{document}
\begin{center}{\Large\bfseries TOAS --- Test 14 (07 Aug 2022)}\end{center}
\begin{multicols}{2}
\raggedcolumns
\partsec{1}{PHYSICS}
\sect{SECTION-I}
\qn{1} 2 mole diatomic gas undergoes cyclic process ABCDA as shown in indicator diagram. Then select correct alternative :-
\qfig{p1_phy_s1_q01_1}
\begin{qoptions}
\opt{A}{Process AB, CD are isobaric and process BC, DA are isochoric.}
\opt{B}{Work done in a cycle is $100R\left\{5-\ln\left(\dfrac{16}{9}\right)\right\}$.}
\opt{C}{Change in entropy of gas and surrounding is zero in a cycle.}
\opt{D}{Change in entropy of gas in process CD is $R\ln\left\{\dfrac{(3)^2(4)^3}{(5)^5}\right\}$.}
\end{qoptions}
\qn{2} At $t=0$ key S is closed then
\qfig{p1_phy_s1_q02_1}
\begin{qoptions}
\opt{A}{Just after closing key $I_2=0$ and $I_1=I=0.1$ A.}
\opt{B}{Long after closing key $I=0.15$ A, $I_1=0.05$ A \& $I_2=0.1$ A.}
\opt{C}{Time constant of growth circuit is 0.1 s.}
\opt{D}{Time constant of growth circuit is 0.075 s.}
\end{qoptions}
\qn{3} Potentiometer wire AB has length 10m and resistance $50\Omega$. When only key $S_1$ is closed balanced length is 2m, when only key $S_2$ is closed balance length is 3m but when both keys are closed then balanced length is 2.5 m. Select correct alternatives.
\qfig{p1_phy_s1_q03_1}
\begin{qoptions}
\opt{A}{$\varepsilon_1=10$V}
\opt{B}{$\varepsilon_2=15$V}
\opt{C}{$r_1=r_2$}
\opt{D}{Potential gradient on wire $=5$V/m}
\end{qoptions}
\qn{4} A $+Q$ charged is fixed on a smooth horizontal table. A particle of mass $m$ and charge $q$ is placed at distance $r$ from $+Q$ charge at given speed $v$. Uniform magnetic field $B$ is inside plane of paper as shown. ($q/m=10^{-2}$ C/kg, $Q=10^{-6}$C, $r=10$m, $B=100$ T)
\qfig{p1_phy_s1_q04_1}
\begin{qoptions}
\opt{A}{For $v=1$ m/s, particle will move on a circular path of radius 10m with constant speed.}
\opt{B}{For $v=0.9$ m/s particle will moving along a line with constant speed.}
\opt{C}{For $v=5$ m/s, when particle reaches a distance 5m from $+Q$ it's speed becomes $\sqrt7$ m/s.}
\opt{D}{For $v=9$ m/s, particle will move on a circular path of radius 10m with constant speed.}
\end{qoptions}
\qn{5} A cylinder of radius 10cm is filled with ideal gas and is filled with a piston made by connecting two solid cylinder of same thickness 5cm but different radius 10cm and 20cm as shown in figure. Initially when 3cm of piston is above water, only 2cm length of piston is inside cylinder, atmospheric pressure is $10^5$ Pa and mass of piston is 5kg. Select correct alternatives.
\qfig{p1_phy_s1_q05_1}
\begin{qoptions}
\opt{A}{Initial pressure of ideal gas is less than $10^5$ Pa.}
\opt{B}{Initial pressure of ideal gas is more than $10^5$ Pa.}
\opt{C}{If water level increased by 2cm pressure of ideal gas will remain same.}
\opt{D}{If water level is increased by 2 cm pressure of gas will decrease.}
\end{qoptions}
\qn{6} When a ray parallel to base of right angled prism incident as shown in figure after two total internal reflection it emerge parallel to incident ray then :
\qfig{p1_phy_s1_q06_1}
\begin{qoptions}
\opt{A}{Refractive index of prism is $\dfrac{\cos\theta}{\cos3\theta}$.}
\opt{B}{Refractive index of prism is $\dfrac{\sin3\theta}{\sin\theta}$.}
\opt{C}{$\dfrac{h_2}{h_1}=\dfrac{\sin\theta}{\sin3\theta}$}
\opt{D}{$\dfrac{h_2}{h_1}=\dfrac{\cos3\theta}{\cos\theta}$}
\end{qoptions}
\qfullbegin
\qn{7} A conducting ring of radius $R$ is moving with constant velocity $v$, at $t=0$ it start entering region containing uniform magnetic field $B$ inside plane of paper then select correct alternatives.
\qfig{p1_phy_s1_q07_1}
\qfullsplit
\begin{qoptions}
\opt{A}{\qfig{p1_phy_s1_q07_A}}
\opt{B}{\qfig{p1_phy_s1_q07_B}}
\opt{C}{\qfig{p1_phy_s1_q07_C}}
\opt{D}{\qfig{p1_phy_s1_q07_D}\par $P_{\text{loss}}$ : Rate of heat production; $t$ : Time}
\end{qoptions}
\qfullend
\qn{8} Four plates of equal area $A$ and plate seperation $d$ are arranged as shown and filled with different dielectric having dielectric constant $K$. $\left(\dfrac{\epsilon_0A}{d}=1\mu\text{F}\right)$ Select correct alternative.
\qfig{p1_phy_s1_q08_1}
\begin{qoptions}
\opt{A}{Effective capacitance of system is $\dfrac43\mu$F .}
\opt{B}{Charge on $2^{\text{nd}}$ plate is $-10~\mu$C.}
\opt{C}{Charge on $3^{\text{rd}}$ plate is $+\dfrac{40}{3}\mu$C.}
\opt{D}{Energy stored in system is $\dfrac{200}{3}\mu$J.}
\end{qoptions}
\sect{SECTION-II}
\paraq{1}{2}
In a cylindrical region uniform magnetic field (inside plane of paper) is increasing at constant rate $(4\pi+3\sqrt3)\,\dfrac{\text T}{\text s}$. A conducting loop ABCDA is placed as shown. ABC is arc with centre at axis of cylinder. Side AC also has a cell of EMF, $\varepsilon=133/12$ volt, then : ($\pi^2\approx10$)
\qfig{p1_phy_s2_para1_2_1}
\qn{1} Electric current in loop will be (in A)
\qn{2} If potential difference between points C and D is $x$ volt then $\dfrac{x}{100}$ is (take $4\pi\sqrt3\approx22$ )
\paraq{3}{4}
Two parallel conducting rails are inclined at angle $\theta$ from horizontal. Their upper ends are connected by an inductor coil of inductance $L$. A metallic rod of mass $m$ starts sliding down. The rail and a uniform magnetic field $B$ acts as shown. ($\theta=30^\circ$, $L=2$H, $m=2$kg, $\ell=1$m and $B=2$T, $g=10$ m/s$^2$)
\qfig{p1_phy_s2_para3_4_1}
\qn{3} Time (in s) at which rod will lose contact with rail :
\qn{4} Velocity of rod (in m/s), just before it lose contact with rails :
\paraq{5}{6}
A hemisphere of radius $R$ is removed from a cylinder of radius $R$ and length $3R$. This cylinder is made of material having density $\sigma=2\rho$ and it is submerged in a fluid of density $\rho$ as shown in figure with help of a string.
\qfig{p1_phy_s2_para5_6_1}
\qn{5} If tension in string is $x\rho(\pi R^3)y$ then value of $x$ is -
\qn{6} If force applied by fluid on hemispherical surface is $y\pi R^3\rho g$ then value of $y$ is -
\sect{SECTION-III}
\qn{1} At $t=0$ key is closed then find maximum current through coil (in mA).
\qfig{p1_phy_s3_q01_1}
\qn{2} In an experiment an oil drop fall with constant speed $v$ in a viscous fluid. When a downward uniform electric field is applied drop fall with constant speed $2v$. When uniform electric field $\alpha E$ is applied in upward direction than drop rises with constant speed $v$. Find $\alpha$.
\qn{3} The self inductance of two coaxial long cylindrical shells of radius $R$ and $2R$, carrying equal current along their length but in opposite direction is $L=\dfrac{\mu_0\ell}{\alpha\pi}\ln(\beta)$, where $\ell$ is length of cylinder then $\alpha+\beta$ is
\qn{4} Figure shows a square loop 10cm on each side in x--y plane with its centre at the origin. An infinite long wire at $z=12$ cm above y--axis.
\qfig{p1_phy_s3_q04_1}
If magnetic torque on loop is expressed as $x\times10^{-7}$ Nm fill value of $x/90$.
\qn{5} Consider a rectangular block of wood moving with velocity $v_0=5$m/s in a gas at temperature $27^\circ$C and mass density 2 kg/m$^3$. Velocity of block is perpendicular to it's face having area $A=1$ cm$^2$. Find drag force (in N) on the block of molar mass of gas is 10 g/mol $\left(R=\dfrac{25}{3}\dfrac{\text J}{\text{mol K}}\right)$.
\qn{6} Semi-circular conducting wire of radius $R=1$m, has it's one end pivoted at origin and it is rotating about z-axis with constant angular velocity $\omega=10$ rad/s. A current carrying long wire is along z-axis carrying current $I=2$A in $+z$ direction. Find potential difference across ends of wire.
\qfig{p1_phy_s3_q06_1}
\partsec{2}{CHEMISTRY}
\sect{SECTION-I}
\qn{1} A one litre solution contains 0.04 M Na$_2$HPO$_4$ and 0.02 M NaH$_2$PO$_4$. For H$_3$PO$_4$ : $\text{p}K_{a_1}=2.14$ ; $\text{p}K_{a_2}=6.84$ ; $\text{p}K_{a_3}=12.3$. The correct statement (s) is /are :\par [Assume volume of solution is constant]\par (use : $\log2=0.3$, $\log5=0.7$) :
\begin{qoptions}
\opt{A}{pH of solution is 7.14}
\opt{B}{On adding 0.02 moles of HCl to initial solution pH increases by 0.6 units}
\opt{C}{On adding 0.02 moles of HCl to initial solution pH decreases by 0.6 units}
\opt{D}{On adding 0.01 mol NaOH to initial solution pH of solution becomes 7.54}
\end{qoptions}
\qn{2} Which of the following option(s) is/are correct ?
\begin{qoptions}
\opt{A}{Angular functions of $2p_x$ and $3p_x$ orbitals in H-atom are identical}
\opt{B}{Radial functions of orbitals in H-atom mainly determine their energy and size}
\opt{C}{Net magnetic field of an atom is due to combined effect of orbital motion and spin of all its electrons in various shells}
\opt{D}{Radial functions of $2p_x$ \& 2s orbitals are identical}
\end{qoptions}
\qn{3} Which of the following match is/are correct?
\begin{qoptions}
\opt{A}{Extraction Method - Self reduction\\ Metal(s) - Hg, Fe}
\opt{B}{Extraction Method - Electrolytic reduction\\ Metal(s) - Na, Al}
\opt{C}{Extraction Method - Hydrometallurgy\\ Metal(s) - Ag, Au}
\opt{D}{Extraction Method - Alumino-thermite reduction\\ Metal(s) - Mn, Cr}
\end{qoptions}
\qn{4} The correct statement(s) related to the metal extraction processes is(are)
\begin{qoptions}
\opt{A}{A mixture of PbS and PbO undergoes self-reduction to produce Pb and SO$_2$.}
\opt{B}{In the extraction process of copper from copper pyrites, silica is added to produce copper silicate.}
\opt{C}{Partial oxidation of sulphide ore of copper by roasting, followed by self-reduction produces blister copper.}
\opt{D}{In cyanide process, zinc powder is utilized to precipitate gold from Na[Au(CN)$_2$].}
\end{qoptions}
\qn{5} Which of the following compounds will produce a compound with offensive smell when treated with CHCl$_3$/KOH?
\begin{qoptions}
\opt{A}{\qfigi{p1_chem_s1_q05_A}}
\opt{B}{\qfigi{p1_chem_s1_q05_B}}
\opt{C}{\qfigi{p1_chem_s1_q05_C}}
\opt{D}{\qfigi{p1_chem_s1_q05_D}}
\end{qoptions}
\qn{6} Which of the following will be the correct representation for products P$_1$, P$_2$ :\par
\noindent CH$_3$--C(=O)--$^{18}$O--Et $\xrightarrow[\text{H}_2\text{O}]{\text{H}^{\oplus}}$ P$_1$ + P$_2$
\begin{qoptions}
\opt{A}{CH$_3$--C(=O)--OH}
\opt{B}{CH$_3$--C(=$^{18}$O)--OH}
\opt{C}{Et--$^{18}$OH}
\opt{D}{Et--OH}
\end{qoptions}
\qn{7}
\qfig{p1_chem_s1_q07_1}
\begin{qoptions}
\opt{A}{Me--C(=O)--OH}
\opt{B}{Et--C(=O)--OH}
\opt{C}{Me--NH$_2$}
\opt{D}{Et--NH$_2$}
\end{qoptions}
\qn{8} $\text{MnO}_4^- + \text{Mn}^{2+}\xrightarrow[\text{H}^+]{\text{catalyst}}$ redox change\par
Select correct for above reaction.
\begin{qoptions}
\opt{A}{It is disproportionation reaction}
\opt{B}{ZnO acts as catalyst}
\opt{C}{ZnSO$_4$ acts as catalyst}
\opt{D}{$\text{MnO}_4^{2-}$ is one of the product}
\end{qoptions}
\sect{SECTION-II}
\paraq{1}{2}
When 100 mL of 1.0 M HCl was mixed with 100 mL of 1.0 M NaOH in an insulated beaker at constant pressure, a temperature increase of 5.7 $^\circ$C was measured for the beaker and its contents.\par
(Expt-1). Because the enthalpy of neutralisation of a strong acid with a strong base is a constant ($-57.0$ kJmol$^{-1}$), this experiment could be used to measure the calorimeter constant. In a second experiment (Expt-2), 100 mL of 2.0 M acetic acid ($K_a=2.0\times10^{-5}$) was mixed with 100 mL of 1.5 M NaOH (under identical conditions to (Expt-1)) where a temperature rise of 8.25 $^\circ$C was measured.\par
(Consider heat capacity of all solutions as 4.2 Jg$^{-1}$K$^{-1}$ and density of all solutions as 1.0 g mL$^{-1}$)
\qn{1} Enthalpy of dissociation (in kJ mol$^{-1}$) of acetic acid obtained from the Expt-2 is
\qn{2} The pH of the solution after Expt-2
\paraq{3}{4}
\qfig{p1_chem_s2_q02_1}
\qn{3} Difference in average oxidation number of S-atom present in anionic part of X and Y.\\ (answer the magnitude value)
\qn{4} When Y is made to react with KI$_{3(\text{aq.})}$ solution then the average oxidation number of S-atom present in product will be.
\paraq{5}{6}
In a polyhydric alcohol, the total number of hydroxyl groups are estimated by the following methods:\par
(i) A known mass of the acetyl derivative of polyhdric alcohol is taken. It is mixed with an excess volume of a standard alkali solution. Both are refluxed. Alkali is used in the hydrolysis of acetyl derivative. The unused alkali is measured by doing titration with a standard acid solution and total number of hydroxyl goups (n) can thus be calculated stoichiometrically.\par
{\footnotesize\noindent A(OH)$_n$ $\xrightarrow{n\text{CH}_3\text{COCl}}$ A(OCOCH$_3$)$_n$ $\xrightarrow{n\text{KOH}}$ A(OH)$_n$ + $n$CH$_3$COOK\par}
{\footnotesize\noindent Polyhydric alcohol \hfill Acetyl derivative\par}
(ii) Total number of-OH groups in vicinal position can be calculated by periodic acid or lead tetra acetate oxidation.
\qfig{p1_chem_s2_para5_6_2}
\qn{5} 0.436 g of acetyl derivative of a polyhydric alcohol (molecular mass = 92) requires 0.336 g KOH for hydrolysis. The total numbers of hydroxyl group(s) in the alcohol are
\qn{6} A compound with molecular formula C$_4$H$_{10}$O$_3$ is converted by the action of acetyl chloride to a compound(s) with molecular weight 190. How many number of acetyl group present in compound (B)
\sect{SECTION-III}
\qn{1} Number of crystal system having at least 2 interfacial angles equal, is $x$\\
Number of crystal system having none of the interfacial angles equal, is $y$\\
Number of crystal system having none of the axial length to be equal, is $z$\\
Find the value of $x-y+z$.
\qn{2} Two moles of an ideal monoatomic gas is taken from state A to state B through a process AB in which the pressure and the temperature of the gas vary such that $PT=$ constant. The process can be represented on a P--T graph as follows ($T_1=300$ K) :
\qfig{p1_chem_s3_q02_1}
The magnitude of heat involved during process AB is (Give your answer after dividing by $300R$) :
\qn{3} When the formation of one mole HBr(g) takes place reversibly from gaseous H$_2$ and liq. Br$_2$ at 300K and 1 atm increase in entropy of the surrounding is 170 J/K. What is change in entropy (J/K) of the system if reaction $2\text{HBr(g)}\to\text{H}_2\text{(g)}+\text{Br}_2\text{(l)}$ is carried out reversibly at 300K in rigid container. [Given : R =8.3 J/mol-K]\\
\textbf{\textit{Fill your answer as sum of digits (excluding decimal places) till you get the single digit answer.}}
\qn{4} Find number of process or ore given in Column-I are correctly represented with Column-II. You have to consider P for P, Q for Q, R for R etc.\par
\smallskip
{\footnotesize\setlength{\tabcolsep}{2pt}\noindent\begin{tabular}{|L{0.55cm}L{2.7cm}|L{0.55cm}L{2.7cm}|}
\hline
\multicolumn{2}{|c|}{(I)\ Process/Ore} & \multicolumn{2}{c|}{(II)\ Related to}\\ \hline
(P) & Matte formation & (P) & Cu\\ \hline
(Q) & Cynide leaching & (Q) & Au\\ \hline
(R) & Monds process & (R) & Reaction of metal $+$ H$_2$\\ \hline
(S) & Fused oxide electrolysis & (S) & Al\\ \hline
(T) & Chromite ore & (T) & Fe $+$ Cr\\ \hline
(U) & Van Arkel process & (U) & Zr\\ \hline
(V) & Sulphuric acid leaching & (V) & Low grade Cu ore\\ \hline
(W) & Separation of FeWO$_4$ & (W) & Extraction of Pb\\ \hline
\end{tabular}\par}
\qn{5} Number of compounds which on acidic hydrolysis gives carboxylic acid as one of the product :\par
(i) CH$_3$CN \quad (ii) CH$_3$NC \quad (iii) CH$_3$CONH$_2$ \quad (iv) CH$_3$COCl\par
(v) (CH$_3$CO)$_2$O \quad (vi) CH$_3$COOCH$_3$ \quad (vii) CH$_3$CHO \quad (viii) CH$_3$CH$_2$OH\par
(ix) CH$_3$COCH$_3$ \quad (x) CH$_3$OCH$_3$
\qfullbegin
\qn{6} Consider following carboxylic acids :
\qfig{p1_chem_s3_q06_1}
\qfullsplit
\qfig{p1_chem_s3_q06_2}
\qfig{p1_chem_s3_q06_3}
\qfig{p1_chem_s3_q06_4}
\qfig{p1_chem_s3_q06_5}
If :\par
(i) The number of given carboxylic acids which exhibit stereoisomerism \& on heating makes cyclic anhydride is $x$ ;\par
(ii) The number of given carboxylic acids which does not form any cyclic anhydride is $y$.\par
Report your answer as $x+y$?
\qfullend
\partsec{3}{MATHEMATICS}
\sect{SECTION-I}
\qn{1} $\displaystyle\int\dfrac{dx}{\sqrt2+\sin x+\cos x}$ is equal to -
\begin{qoptions}
\opt{A}{$\dfrac1{\sqrt2}\left(\sec\left(x+\dfrac\pi4\right)-\tan\left(x+\dfrac\pi4\right)\right)+C$}
\opt{B}{$\dfrac1{\sqrt2}\left(\sec\left(x-\dfrac\pi4\right)-\tan\left(x-\dfrac\pi4\right)\right)+C$}
\opt{C}{$\dfrac1{\sqrt2}\left(\tan\left(x+\dfrac\pi4\right)-\sec\left(x+\dfrac\pi4\right)\right)+C$}
\opt{D}{$\dfrac1{\sqrt2}\left(\tan\left(x-\dfrac\pi4\right)-\sec\left(x-\dfrac\pi4\right)\right)+C$}
\end{qoptions}
\qn{2} A function $f:R\to R$ satisfies the differential equation $2xy+(1+x^2)y'=1$, (where $y'$ denotes derivative of $y$ with respect to $x$) where $f(0)=0$, then
\begin{qoptions}
\opt{A}{$f(x)$ is odd}
\opt{B}{$f(x)=1$ has no real solution}
\opt{C}{$f(x)=\dfrac12$ has exactly two solution}
\opt{D}{$f(x)=-1$ has no real solution}
\end{qoptions}
\qn{3} Which of the following options is/are true ?
\begin{qoptions}
\opt{A}{for $\alpha>1$, $(1+x)^\alpha\ge1+\alpha x\ \forall\ x>-1$}
\opt{B}{for $\alpha>1$, $(1+x)^\alpha<1+\alpha x\ \forall\ x>-1$}
\opt{C}{for $x>0$, $e^{\ln x}\le x$}
\opt{D}{for $x>0$, $e^{\ln x}>x$}
\end{qoptions}
\qn{4} Area enclosed by parabola $ay=3(a^2-x^2)$ and x-axis is 64, then value of `$a$' can be -
\begin{qoptions}
\opt{A}{4}
\opt{B}{2}
\opt{C}{$-4$}
\opt{D}{$-2$}
\end{qoptions}
\qn{5} Let $A$ be the area bounded by the curve $f(x)=\dfrac{2}{1+x^4}$, x-axis, y-axis and ordinate $x=1$, then -
\begin{qoptions}
\opt{A}{$A>\dfrac\pi2$}
\opt{B}{$A>2\ln2$}
\opt{C}{$0<A<\dfrac\pi2$}
\opt{D}{$\dfrac\pi2<A<2\ln2$}
\end{qoptions}
\qn{6} Let $C$ be a curve such that normal at any point $P$ on it meets x-axis and y-axis at $A$ and $B$ respectively. If $BP:PA=1:2$ (internally) and curve passes through $(0,4)$, then -
\begin{qoptions}
\opt{A}{The curve `$C$' passes through $(\sqrt{10},-6)$}
\opt{B}{Equation of tangent at $(4,4\sqrt3)$ is $\sqrt3x-2y+4\sqrt3=0$}
\opt{C}{Equation of tangent at $(4,4\sqrt3)$ is $2x-\sqrt3y+4=0$}
\opt{D}{If point $P$ is $(8,12)$, then point $A$ and point $B$ are $(24,0)$ and $(0,18)$ respectively.}
\end{qoptions}
\qn{7} Let $S$ be the area bounded by $y=e^{|\cos4x|}$, $x=0$, $y=0$ and $x=\pi$. Then
\begin{qoptions}
\opt{A}{$S=2\displaystyle\int_0^{\pi/2}e^{\sin t}\,dt$}
\opt{B}{$S>\dfrac\pi2$}
\opt{C}{$S<2\left(e^{\frac\pi2}-1\right)$}
\opt{D}{$S\le1$}
\end{qoptions}
\qn{8} The curve(s) which satisfy the differential equation $x\,dy=y\,dx+\sin\left(\dfrac1x\right)dx$ is/are -
\begin{qoptions}
\opt{A}{$y=x\left(\cos\left(\dfrac1x\right)+3\right)$}
\opt{B}{$y=2x\left(2\cos^2\dfrac1{2x}+\sin^2\dfrac1{2x}\right)$}
\opt{C}{$y=x^2\left(\cos\left(\dfrac1x\right)+3\right)$}
\opt{D}{$y=2x\left(2\cos^2\dfrac1{2x}-\sin^2\dfrac1{2x}\right)$}
\end{qoptions}
\sect{SECTION-II}
\paraq{1}{2}
Let $g(x)=\dfrac{x(x-1)}{(x^2+1)}e^{-x}$ and function $y=f(x)$ satisfy the differential equation $dy-y\,dx=e^xg'(x)\,dx$ also $f(0)=0$
\qn{1} $f'(1)$ is equal to -
\qn{2} If area bounded by $f(x)$ and x-axis in IV quadrant is $\dfrac\pi a-\ln\left(\dfrac e{\sqrt2}\right)$, then $a$ is equal to -
\paraq{3}{4}
Consider a function $f(x)=\dfrac{x^3}{3}-\dfrac{x^2}{2}+x+\dfrac1{12}$.
\qn{3} The value of $\displaystyle\sum_{r=1}^{99}f\left(\dfrac r{100}\right)$ is -
\qn{4} The value of integral $\displaystyle\int_{1/7}^{6/7}f(f(x))\,dx$ is -
\paraq{5}{6}
Let $f:R\to R$ and $g:R\to R$ such that\par
$f(x)=\max\{|x|,\ |x-1|,|x-2|\ .....|x-2018|\}$\par
$g(x)=\min\{|x|,|x-1|,|x-2|\ ....|x-2018|\}$\ \ then :
\qn{5} The value of $\displaystyle\int_0^{2018}f(x)\,dx$ is $3(a)^2$ then value of $a$ is :
\qn{6} The value of $\dfrac{\displaystyle\int_0^{2018}f(x)\,dx}{\displaystyle\int_0^{2018}g(x)\,dx}$ is
\sect{SECTION-III}
\qn{1} Let $f:[0,1]\to(0,\infty)$ be a continuous function such that $\displaystyle\int_0^1f(x)\,dx=1$, then maximum value of $\left(\displaystyle\int_0^1\sqrt[3]{f(x)}\,dx\right)\left(\displaystyle\int_0^1\sqrt[5]{f(x)}\,dx\right)\left(\displaystyle\int_0^1\sqrt[7]{f(x)}\,dx\right)$, is
\qn{2} The value of $\sqrt2+\displaystyle\int_0^1\dfrac{2x^2-2x+1}{\sqrt{x^2-2x+2}}\,dx$ is -
\qn{3} The value of $\displaystyle\int_1^3\left(\sqrt{1+(x-1)^3}+(x^2-1)^{1/3}\right)dx$, is
\qn{4} The curve $y=e^{2x}-18x+15$ crosses the y-axis at $P$ and has point of minima at $Q$. If the area bounded by the curve and chord $PQ$ is $(k\ln3-k+1)$ then $k$ is equal to (use $\ln3=1.1$)
\qn{5} Let $a=\cos^{-1}\cos20$, $b=\cos^{-1}\cos30$, $c=\sin^{-1}\sin(a+b)$ and $5\sec^{-1}\alpha+10\sin^{-1}\beta=10(a+b+c)$ if the value of $\tan^{-1}\alpha+\cos^{-1}(\beta-1)$ is $k$ then value of $\left(\dfrac{12}{\pi}k\right)$ is
\qn{6} If `$e$' denotes the eccentricity of the hyperbola, satisfying the differential equation $2xy\dfrac{dx}{dy}=x^2+y^2$ and passing through the point $(2,1)$, then $(e^2-1)$ is equal to
\end{multicols}
\end{document}
